% Quorums avec N = 3 : W = 2 et R = 2 se recouvrent toujours ; W = 1 et R = 1
% peuvent manquer la réplique qui a la dernière écriture.
\documentclass[border=6pt]{standalone}
\usepackage{sd}
\begin{document}
\begin{tikzpicture}[x=1mm, y=1mm, note/.append style={font=\sffamily\small},
    rep/.style={draw=#1, fill=#1Bg, line width=0.7pt, rounded corners=2pt, minimum width=15mm, minimum height=9mm, font=\sffamily\footnotesize, align=center}]
  \foreach \titre/\dx/\ecr/\lec/\res/\c in {{N = 3, W = 2, R = 2}/0/{1,2}/{2,3}/{R + W = 4 > 3 : la lecture touche\\toujours une réplique à jour}/Dgreen,
                                          {N = 3, W = 1, R = 1}/85/{1}/{3}/{R + W = 2 ≤ 3 : la lecture peut\\ne toucher qu'une réplique ancienne}/Dred} {
    \begin{scope}[xshift=\dx mm]
      \node[pillo=\c] at (35,44) {\titre};
      \foreach \i in {1,2,3} {
        \node[rep=Dslate] (r\i) at ({12+22*(\i-1)},22) {réplique \i\\v = 0};
      }
      \foreach \i in \ecr {
        \node[rep=Dblue] at ({12+22*(\i-1)},22) {réplique \i\\v = 7};
        \draw[arr=Dblue] ({12+22*(\i-1)},35) -- ({12+22*(\i-1)},27.5);
      }
      \node[note, anchor=south, text=Dblue] at (35,35) {écriture v = 7};
      \foreach \i in \lec {
        \draw[arr=Dorange] ({12+22*(\i-1)},16.5) -- ({12+22*(\i-1)},9);
      }
      \node[note, anchor=north, text=Dorange] at (35,8) {lecture};
      \node[note, anchor=north, align=center, text=\c] at (35,1) {\res};
    \end{scope}
  }
\end{tikzpicture}
\end{document}
